\chapter*{Resumen}
\addcontentsline{toc}{chapter}{Resumen}

Anotar manualmente un experimento de nado forzado lleva entre una y dos horas por
video. Un investigador reproduce la grabación, observa la conducta de cada rata y
registra a mano cuánto tiempo pasa nadando, inmóvil o intentando escalar las paredes
del cilindro. Si el protocolo requiere que dos evaluadores lo hagan de forma
independiente para verificar consistencia, ese tiempo se duplica. Con entre 32 y 40
videos por semestre, la anotación consume cientos de horas-persona y produce resultados
que varían según quién anota y cómo interpreta cada transición entre conductas.

Este Trabajo Terminal propone un prototipo de sistema web con visión por computadora y
aprendizaje supervisado que automatice ese proceso para el Laboratorio de Bioquímica
Estructural, Sección de Posgrado, de la Escuela Nacional de Medicina y Homeopatía
(ENMyH-IPN). El uso concreto
es directo: el investigador sube el video al sistema desde el navegador, el sistema lo
analiza y en minutos genera un reporte con los tiempos de nado activo, inmovilidad e
intentos de escape por espécimen y sesión.

El sistema no diagnostica depresión ni produce información de uso clínico. Es una
herramienta de apoyo a la investigación preclínica en farmacología experimental; su
propósito es reemplazar la anotación manual de videos del Forced Swim Test (FST), la
prueba conductual más usada para evaluar el efecto de antidepresivos en ratas antes de
pasar a ensayos en humanos.

El sistema integra un módulo de visión por computadora y aprendizaje supervisado que
clasifica las conductas de interés directamente sobre los videos del laboratorio, y una
plataforma web que permite cargar videos, ejecutar el análisis y obtener reportes
cuantitativos por espécimen y por sesión. La validación compara los resultados del
sistema contra los registros manuales del laboratorio (gold standard) usando el
coeficiente kappa de Cohen para medir concordancia y el error absoluto medio para
cuantificar la diferencia en segundos por conducta. El desarrollo sigue el ciclo de
vida de software de la norma ISO/IEC/IEEE~12207:2017 y los atributos de calidad de la
ISO/IEC~25010:2023.


\noindent\textbf{Palabras clave:} visión por computadora, aprendizaje automático,
prueba de nado forzado, análisis conductual preclínico, sistema web, farmacología
experimental.
